%% TRANSCRIBED (PDF 24-37)

\sechead{2}{\textit{l}-adic representations of number fields}%
  {\textit{l}-ADIC REPRESENTATIONS OF NUMBER FIELDS}

\subsechead{2.1}{Preliminaries}

(For the basic notions concerning number fields, see for instance
Cassels-Fr\"ohlich~\cite{cassels-frohlich}, Lang~\cite{lang-algnum} or
Weil~\cite{weil-bnt}.)  Let $K$ be a number field (i.e. a finite
extension of $\Q$).  Denote by $\Sigma_K$ the set of all finite places
of $K$, i.e., the set of all normalized discrete valuations of $K$
(or, alternatively, the set of prime ideals in the ring $A_K$ of
integers of $K$).  The \emph{residue field} $k_v$ of a place
$v \in \Sigma_K$ is a finite field with $N_v = p_v^{\deg(v)}$
elements, where $p_v = \mathrm{char}(k_v)$ and $\deg(v)$ is the degree
of $k_v$ over $\F_{p_v}$.  The ramification index $e_v$ of $v$ is
$v(p_v)$.

Let $L/K$ be a finite Galois extension with Galois group $G$, and let
$w \in \Sigma_L$.  The subgroup $D_w$ of $G$ consisting of those
$g \in G$ for which $gw = w$ is the \emph{decomposition group} of $w$.
The restriction of $w$ to $K$ is an integral multiple of an element
$v \in \Sigma_K$; by abuse of language, we also say that $v$ is the
restriction of $w$ to $K$, and we write $w \mid v$ (``$w$ divides
$v$'').  Let $L_w$ (resp. $K_v$) be the completion of $L$ (resp. $K$)
with respect to $w$ (resp. $v$).  We have $D_w = \Gal(L_w/K_v)$.  The
group $D_w$ is mapped homomorphically onto the Galois group
$\Gal(\ell_w/k_v)$ of the corresponding residue extension
$\ell_w/k_v$.  The kernel of $G \longrightarrow \Gal(\ell_w/k_v)$ is
the \emph{inertia group} $I_w$ of $w$.
%% NOTE (I-6): the book prints "kernel of G -> Gal(l_w/k_v)";
%% mathematically the source is D_w, but "G" is kept as printed.
The quotient group $D_w/I_w$ is a finite cyclic group generated by the
\emph{Frobenius element} $F_w$; we have $F(\lambda) = \lambda^{N_v}$
for all $\lambda \in \ell_w$.  The valuation $w$ (resp. $v$) is called
\emph{unramified} if $I_w = \{1\}$.  Almost all places of $K$ are
unramified.

If $L$ is an arbitrary algebraic extension of $\Q$, one defines
$\Sigma_L$ to be the projective limit of the sets $\Sigma_{L_a}$, where
$L_a$ ranges over the finite sub-extensions of $L/\Q$.  Then, if $L/K$
is an arbitrary Galois extension of the number field $K$, and
$w \in \Sigma_L$, one defines $D_w$, $I_w$, $F_w$ as before.  If $v$
is an unramified place of $K$, and $w$ is a place of $L$ extending
$v$, we denote by $F_v$ the conjugacy class of $F_w$ in
$G = \Gal(L/K)$.

\begin{dfn}
Let $\rho : \Gal(\bar{K}/K) \longrightarrow \Aut(V)$ be an
\textit{l}-adic representation of $K$, and let $v \in \Sigma_K$.  We
say that $\rho$ is unramified at $v$ if $\rho(I_w) = \{1\}$ for any
valuation $w$ of $\bar{K}$ extending $v$.
\end{dfn}

If the representation $\rho$ is unramified at $v$, then the
restriction of $\rho$ to $D_w$ factors through $D_w/I_w$ for any
$w \mid v$; hence $\rho(F_w) \in \Aut(V)$ is defined; we call
$\rho(F_w)$ the \emph{Frobenius} of $w$ in the representation $\rho$,
and we denote it by $F_{w,\rho}$.  The conjugacy class of $F_{w,\rho}$
in $\Aut(V)$ depends only on $v$; it is denoted by $F_{v,\rho}$.  If
$L/K$ is the extension of $K$ corresponding to $H = \Ker(\rho)$, then
$\rho$ is unramified at $v$ if and only if $v$ is unramified in $L/K$.

\subsechead{2.2}{\v{C}ebotarev's density theorem}

Let $P$ be a subset of $\Sigma_K$.  For each integer $n$, let
$a_n(P)$ be the number of $v \in P$ such that $N_v \leq n$.  If $a$ is
a real number, one says that $P$ \emph{has density} $a$ if
\[
\lim.\ \frac{a_n(P)}{a_n(\Sigma_K)} = a
\qquad\text{when}\quad n \longrightarrow \infty.
\]
Note that $a_n(\Sigma_K) \sim n/\log(n)$, by the prime number theorem
(cf.~Appendix, or \cite{lang-algnum}, chap.~VIII), so that the above
relation may be rewritten:
\[
a_n(P) = a.n/\log(n) + o(n/\log(n)).
\]

\begin{examples}
A finite set has density $0$.  The set of $v \in \Sigma_K$ of degree
$1$ (i.e. such that $N_v$ is prime) has density $1$.  The set of
ordinary prime numbers whose first digit (in the decimal system, say)
is $1$ has no density.
\end{examples}

We can now state \v{C}ebotarev's density theorem:

\begin{theorem}
Let $L$ be a finite Galois extension of the number field $K$, with
Galois group $G$.  Let $X$ be a subset of $G$, stable by conjugation.
Let $P_X$ be the set of places $v \in \Sigma_K$, unramified in $L$,
such that the Frobenius class $F_v$ is contained in $X$.  Then $P_X$
has density equal to $\mathrm{Card}(X)/\mathrm{Card}(G)$.
\end{theorem}

For the proof, see \cite{cebotarev}, \cite{artin-collected}, or the
Appendix.

\begin{corollary}\label{cor:I.2.2.1}
For every $g \in G$, there exist infinitely many unramified places
$w \in \Sigma_L$ such that $F_w = g$.
\end{corollary}

For infinite extensions, we have:

\begin{corollary}\label{cor:I.2.2.2}
Let $L$ be a Galois extension of $K$, which is unramified outside a
finite set $S$.
\begin{enumerate}[label=\alph*),leftmargin=2.2em,itemsep=0.15ex,topsep=0.4ex]
\item The Frobenius elements of the unramified places of $L$ are dense
in $\Gal(L/K)$.
\item Let $X$ be a subset of $\Gal(L/K)$, stable by conjugation.
Assume that the boundary of $X$ has measure zero with respect to the
Haar measure $\mu$ of $X$, and normalize $\mu$ such that its total
mass is $1$.  Then the set of places $v \notin S$ such that
$F_v \subset X$ has a density equal to $\mu(X)$.
%% NOTE (I-8): "Haar measure mu of X" is as printed (one would expect
%% "of Gal(L/K)").
\end{enumerate}
\end{corollary}

Assertion (b) follows from the theorem, by writing $L$ as an
increasing union of finite Galois extensions and passing to the limit
(one may also use Prop.~1 of the Appendix).  Assertion (a) follows
from (b) applied to a suitable neighborhood of a given class of
$\Gal(L/K)$.

\runhead{Exercise}

Let $G$ be an \textit{l}-adic Lie group and let $X$ be an analytic
subset of $G$ (i.e. a set defined by the vanishing of a family of
analytic functions on $G$).  Show that the boundary of $X$ has measure
zero with respect to the Haar measure of $G$.

\subsechead{2.3}{Rational \textit{l}-adic representations}

Let $\rho$ be an \textit{l}-adic representation of the number field
$K$.  If $v \in \Sigma_K$, and if $v$ is unramified with respect to
$\rho$, we let $P_{v,\rho}(T)$ denote the polynomial
$\det(1 - F_{v,\rho}T)$.

\begin{dfn}
The \textit{l}-adic representation $\rho$ is said to be rational
(resp. integral) if there exists a finite subset $S$ of $\Sigma_K$
such that
\begin{alist}
\item Any element of $\Sigma_K - S$ is unramified with respect to
$\rho$.
\item If $v \notin S$, the coefficients of $P_{v,\rho}(T)$ belong to
$\Q$ (resp. to $\Z$).
\end{alist}
\end{dfn}

\runhead{Remark}

Let $K'/K$ be a finite extension.  An \textit{l}-adic representation
$\rho$ of $K$ defines (by restriction) an \textit{l}-adic
representation $\rho_{/K'}$ of $K'$.  If $\rho$ is rational
(resp. integral), then the same is true for $\rho_{/K'}$; this follows
from the fact that the Frobenius elements relative to $K'$ are powers
of those relative to $K$.

\begin{examples}
The \textit{l}-adic representations of $K$ given in examples 1, 2, 3
of \S\ref{sec:I.1.2} are rational (even integral) representations.  In
example 1, one can take for $S$ the set $S_\ell$ of elements $v$ of
$\Sigma_K$ with $p_v = \ell$; the corresponding Frobenius is $N_v$,
viewed as an element of $U_\ell$.  In examples 2, 3, one can take for
$S$ the union of $S_\ell$ and the set $S_A$ where $A$ has ``bad
reduction''; the fact that the corresponding Frobenius has an integral
characteristic polynomial (which is independent of $\ell$) is a
consequence of Weil's results on endomorphisms of abelian varieties
(cf.~\cite{weil-abelian} and \cite{lang-abelian}, chap.~VII).  The
rationality of the cohomology representations is a well-known open
question.
\end{examples}

\begin{dfn}
Let $\ell'$ be a prime, $\rho'$ an $\ell'$-adic representation of $K$,
and assume that $\rho$, $\rho'$ are rational.  Then $\rho$, $\rho'$
are said to be compatible if there exists a finite subset $S$ of
$\Sigma_K$ such that $\rho$ and $\rho'$ are unramified outside of $S$
and $P_{v,\rho}(T) = P_{v,\rho'}(T)$ for $v \in \Sigma_K - S$.
\end{dfn}

(In other words, the characteristic polynomials of the Frobenius
elements are the same for $\rho$ and $\rho'$, at least for almost all
$v$'s.)

If $\rho : \Gal(\bar{K}/K) \longrightarrow \Aut(V)$ is a rational
\textit{l}-adic representation of $K$, then $V$ has a composition
series
\[
V = V_o \supset V_1 \supset \ldots \supset V_q = 0
\]
of $\rho$-invariant subspaces with $V_i/V_{i+1}$ ($0 \leq i \leq q-1$)
\emph{simple} (i.e. \emph{irreducible}).  The \textit{l}-adic
representation $\rho'$ of $K$ defined by
$V' = \sum\limits_{i=0}^{q-1} V_i/V_{i+1}$ is semi-simple, rational,
and compatible with $\rho$; it is the ``semi-simplification'' of $V$.

\begin{theorem}\label{thm:I.2.3.unnum}
Let $\rho$ be a rational \textit{l}-adic representation of $K$, and
let $\ell'$ be a prime.  Then there exists at most one (up to
isomorphism) $\ell'$-adic rational representation $\rho'$ of $K$ which
is semi-simple and compatible with $\rho$.
\end{theorem}

(Hence there exists a unique (up to isomorphism) rational, semi-simple
\textit{l}-adic representation compatible with $\rho$.)

\begin{proof}
Let $\rho'_1$, $\rho'_2$ be semi-simple $\ell'$-adic representations
of $K$ which are rational and compatible with $\rho$.  We first prove
that $\mathrm{Tr}(\rho'_1(g)) = \mathrm{Tr}(\rho'_2(g))$ for all
$g \in G$.  Let $H = G/(\Ker(\rho'_1) \cap \Ker(\rho'_2))$; the
representations $\rho'_1$, $\rho'_2$ may be regarded as
representations of $H$, and it suffices to show that
$\mathrm{Tr}(\rho'_1(h)) = \mathrm{Tr}(\rho'_2(h))$ for all $h \in H$.
Let $M \subset \bar{K}$ be the fixed field of $H$.  Then by the
compatibility of $\rho'_1$, $\rho'_2$ there is a finite subset $S$ of
$\Sigma_K$ such that for all $v \in \Sigma_K - S$, $w \in \Sigma_M$,
$w \mid v$, we have
$\mathrm{Tr}(\rho'_1(F_w)) = \mathrm{Tr}(\rho'_2(F_w))$.  But, by
cor.~\ref{cor:I.2.2.2} to \v{C}ebotarev's theorem
(cf.~\S\ref{sec:I.2.2}) the $F_w$ are dense in $H$.  Hence
$\mathrm{Tr}(\rho'_1(h)) = \mathrm{Tr}(\rho'_2(h))$ for all $h \in H$
since $\mathrm{Tr}\circ\rho'_1$, $\mathrm{Tr}\circ\rho'_2$ are
continuous.

The theorem now follows from the following result applied to the
group ring $\Lambda = \Q_\ell[H]$.

\begin{lemma}
Let $k$ be a field of characteristic zero, let $\Lambda$ be a
$k$-algebra, and let $\rho_1$, $\rho_2$ be two finite-dimensional
linear representations of $\Lambda$.  If $\rho_1$, $\rho_2$ are
semi-simple and have the same \emph{trace}
($\mathrm{Tr}\circ\rho_1 = \mathrm{Tr}\circ\rho_2$), then they are
isomorphic.
\end{lemma}

For the proof see Bourbaki, \emph{Alg.}, ch.~8, \S12, n°~1, prop.~3.
\end{proof}

\begin{dfn}
For each prime $\ell$ let $\rho_\ell$ be a rational \textit{l}-adic
representation of $K$.  The system $(\rho_\ell)$ is said to be
compatible if $\rho_\ell$, $\rho_{\ell'}$ are compatible for any two
primes $\ell$, $\ell'$.  The system $(\rho_\ell)$ is said to be
strictly compatible if there exists a finite subset $S$ of $\Sigma_K$
such that:
\begin{alist}
\item Let $S_\ell = \{v \mid p_v = \ell\}$.  Then, for every
$v \notin S \cup S_\ell$, $\rho_\ell$ is unramified at $v$ and
$P_{v,\rho_\ell}(T)$ has rational coefficients.
\item $P_{v,\rho_\ell}(T) = P_{v,\rho_{\ell'}}(T)$ if
$v \notin S \cup S_\ell \cup S_{\ell'}$.
\end{alist}
\end{dfn}

When a system $(\rho_\ell)$ is strictly compatible, there is a
smallest finite set $S$ having properties (a) and (b) above.  We call
it the \emph{exceptional set} of the system.

\begin{examples}
The systems of \textit{l}-adic representations given in examples 1, 2,
3 of \S\ref{sec:I.1.2} are each strictly compatible.  The exceptional
set of the first one is empty.  The exceptional set of example 2
(resp. 3) is the set of places where the elliptic curve (resp. the
abelian variety) has ``bad reduction'', cf.~\cite{serre-tate-good}.
\end{examples}

\runhead{Questions}
\begin{enumerate}[label=\arabic*.,leftmargin=2.2em,itemsep=0.15ex,topsep=0.4ex]
\item Let $\rho$ be a rational \textit{l}-adic representation.  Is it
true that $P_{v,\rho}$ has rational coefficients for \emph{all} $v$
such that $\rho$ is unramified at $v$?

A somewhat similar question is:

Is any compatible system strictly compatible?

\item Can any rational \textit{l}-adic representation be obtained (by
tensor products, direct sums, etc.) from ones coming from
\textit{l}-adic cohomology?

\item Given a rational \textit{l}-adic representation $\rho$ of $K$,
and a prime $\ell'$, does there exist a rational $\ell'$-adic
representation $\rho'$ of $K$ compatible with $\rho$?
$\longrightarrow$ [no: easy counter-examples].

\item Let $\rho$, $\rho'$ be rational $\ell$, $\ell'$-adic
representations of $K$ which are compatible and semi-simple.
\begin{rlist}
\item If $\rho$ is abelian (i.e., if $\Ima(\rho)$ is abelian), is it
true that $\rho'$ is abelian?  (We shall see in
chap.~\ref{chap:III} that this is true at least if $\rho$ is
``locally algebraic''.)  $\longrightarrow$ [yes: this follows from
\cite{henniart}.]
\item Is it true that $\Ima(\rho)$ and $\Ima(\rho')$ are Lie groups of
the same dimension?  More optimistically, is it true that there exists
a Lie algebra $\mathfrak{g}$ over $\Q$ such that
$\mathrm{Lie}(\Ima(\rho)) = \mathfrak{g} \tens_{\Q} \Q_\ell$,
$\mathrm{Lie}(\Ima(\rho')) = \mathfrak{g} \tens_{\Q} \Q_{\ell'}$?
%% NOTE (I-13): the book prints the Lie algebra letter as an
%% underlined "g" (boldface fraktur); transcribed as \mathfrak{g}
%% (cf. \mathfrak{s}_m in ch2_s2.tex).
\end{rlist}

\item Let $X$ be a non-singular projective variety defined over $K$,
and let $i$ be an integer.  Is the $i$-th cohomology representation
$H^i_\ell(X_s)$ semi-simple?  Does its Lie algebra contain the
homotheties if $i \geq 1$?  (When $i = 1$, an affirmative answer to
either one of these questions would imply a positive solution for the
``congruence subgroup problem'' on abelian varieties,
cf.~\cite{serre-congr}, \S3.)  $\longrightarrow$ [yes for $i=1$: see
\cite{bogomolov} and also \cite{serre-congr2}.]
\end{enumerate}

\runhead{Remark}

The concept of an \textit{l}-adic representation can be generalized by
replacing the prime $\ell$ by a place $\lambda$ of a number field $E$.
A $\lambda$-adic representation is then a continuous homomorphism
$\Gal(K_s/K) \longrightarrow \Aut(V)$, where $V$ is a
finite-dimensional vector space over the local field $E_\lambda$.  The
concepts of rational $\lambda$-adic representation, compatible
representations, etc., can be defined in a way similar to the
\textit{l}-adic case.

\begin{exercises}
\begin{nlist}
\item Let $\rho$ and $\rho'$ be two rational, semi-simple, compatible
representations.  Show that, if $\Ima(\rho)$ is finite, the same is
true for $\Ima(\rho')$ and that $\Ker(\rho) = \Ker(\rho')$.  (Apply
exer.~3 of \S\ref{sec:I.1.1} to $\rho'$ and to $U = \Ker(\rho)$.)

Generalize this to $\lambda$-adic representations (with respect to a
number field $E$).

\item Let $\rho$ (resp. $\rho'$) be a rational \textit{l}-adic
(resp. $\ell'$-adic) representation of $K$, of degree $n$.  Assume
$\rho$ and $\rho'$ are compatible.  If $s \in G = \Gal(\bar{K}/K)$,
let $\sigma_i(s)$ (resp. $\sigma'_i(s)$) be the $i$-th coefficient of
the characteristic polynomial of $\rho(s)$ (resp. of $\rho'(s)$).
Let $P(X_o, \ldots, X_n)$ be a polynomial with rational coefficients,
and let $X_P$ (resp. $X'_P$) be the set of $s \in G$ such that
$P(\sigma_o(s), \ldots, \sigma_n(s)) = 0$ (resp.
$P(\sigma'_o(s), \ldots, \sigma'_n(s)) = 0$).
\begin{enumerate}[label=\alph*),leftmargin=2.2em,itemsep=0.15ex,topsep=0.4ex]
\item Show that the boundaries of $X_P$ and $X'_P$ have measure zero
for the Haar measure $\mu$ of $G$ (use Exer.~of \S\ref{sec:I.2.2}).
\item Assume that $\mu$ is normalized, i.e. $\mu(G) = 1$.  Let $T_P$
be the set of $v \in \Sigma_K$ at which $\rho$ is unramified, and for
which the coefficients $\sigma_o, \ldots, \sigma_n$ of the
characteristic polynomial of $F_{v,\rho}$ satisfy the equation
$P(\sigma_o, \ldots, \sigma_n) = 0$.  Show that $T_P$ has density
equal to $\mu(X_P)$.
\item Show that $\mu(X_P) = \mu(X'_P)$.
\end{enumerate}
\end{nlist}
\end{exercises}

\subsechead{2.4}{Representations with values in a linear algebraic group}

Let $H$ be a linear algebraic group defined over a field $k$.  If $k'$
is a commutative $k$-algebra, let $H(k')$ denote the group of points
of $H$ with values in $k'$.  Let $A$ denote the coordinate ring (or
``affine ring'') of $H$.  An element $f \in A$ is said to be
\emph{central} if $f(xy) = f(yx)$ for any $x, y \in H(k')$ and any
commutative $k$-algebra $k'$.  If $x \in H(k')$, we say that
\emph{the conjugacy class of $x$ in $H$ is rational over} $k$ if
$f(x) \in k$ for any central element $f$ of $A$.

\begin{dfn}
Let $H$ be a linear algebraic group over $\Q$, and let $K$ be a field.
A continuous homomorphism
$\rho : \Gal(K_s/K) \longrightarrow H(\Q_\ell)$ is called an
\textit{l}-adic representation of $K$ with values in $H$.
\end{dfn}

(Note that $H(\Q_\ell)$ is, in a natural way, a topological group and
even an \textit{l}-adic Lie group.)

If $K$ is a number field, one defines in an obvious way what it means
for $\rho$ to be unramified at a place $v \in \Sigma_K$; if $w \mid
v$, one defines the Frobenius element $F_{w,\rho} \in H(\Q_\ell)$ and
its conjugacy class $F_{v,\rho}$.  We say, as before, that $\rho$ is
\emph{rational} if
\begin{alist}
\item there is a finite set $S$ of $\Sigma_K$ such that $\rho$ is
unramified outside $S$,
\item if $v \notin S$, the conjugacy class $F_{v,\rho}$ is rational
over $\Q$.
\end{alist}
Two rational representations $\rho$, $\rho'$ (for primes $\ell$,
$\ell'$) are said to be \emph{compatible} if there exists a finite
subset $S$ of $\Sigma_K$ such that $\rho$ and $\rho'$ are unramified
outside $S$ and such that for any central element $f \in A$ and any
$v \in \Sigma_K - S$ we have
$f(F_{v,\rho}) = f(F_{v,\rho'})$.  One defines in the same way the
notions of \emph{compatible} and \emph{strictly compatible systems} of
rational representations.

\begin{remarks}
\begin{enumerate}[label=\arabic*.,leftmargin=2.2em,itemsep=0.15ex,topsep=0.4ex]
\item If the algebraic group $H$ is abelian, then condition (b) above
means that $F_{v,\rho}$ (which is now an element of $H(\Q_\ell)$) is
\emph{rational over} $\Q$, i.e. belongs to $H(\Q)$.

\item Let $V_o$ be a finite-dimensional vector space over $\Q$, and
let $\mathrm{GL}_{V_o}$ be the linear algebraic group over $\Q$ whose
group of points in any commutative $\Q$-algebra $k$ is
$\Aut(V_o \tens_{\Q} k)$; in particular, if
$V_\ell = V_o \tens_{\Q} \Q_\ell$, then
$\mathrm{GL}_{V_o}(\Q_\ell) = \Aut(V_\ell)$.  If
$\phi : H \longrightarrow \mathrm{GL}_{V_o}$ is a homomorphism of
linear algebraic groups over $\Q$, call $\phi_\ell$ the induced
homomorphism of $H(\Q_\ell)$ into
$\mathrm{GL}_{V_o}(\Q_\ell) = \Aut(V_\ell)$.  If $\rho$ is an
\textit{l}-adic representation of $\Gal(\bar{K}/K)$ into $H(\Q_\ell)$,
one gets by composition a linear \textit{l}-adic representation
$\phi_\ell \circ \rho : \Gal(K_s/K) \longrightarrow \Aut(V_\ell)$.
Using the fact that the coefficients of the characteristic polynomial
are central functions, one sees that $\phi_\ell \circ \rho$ is
\emph{rational} if $\rho$ is \emph{rational} ($K$ a number field).  Of
course, compatible representations in $H$ give compatible linear
representations.  We will use this method of constructing compatible
representations in the case where $H$ is abelian (see
chap.~\ref{chap:II}, \S\ref{sec:II.2.5}).
\end{enumerate}
\end{remarks}

\subsechead{2.5}{\textit{L}-functions attached to rational representations}

Let $K$ be a number field and let $\rho = (\rho_\ell)$ be a strictly
compatible system of rational \textit{l}-adic representations, with
exceptional set $S$.  If $v \notin S$, denote by $P_{v,\rho}(T)$ the
rational polynomial $\det(1 - F_{v,\rho_\ell}T)$, for any
$\ell \neq p_v$; by assumption, this polynomial does not depend on the
choice of $\ell$.  Let $s$ be a complex number.  One has:
\begin{align*}
P_{v,\rho}(N_v^{-s}) &= \det(1 - F_{v,\rho}/(N_v)^s) \\
&= \prod_i (1 - \lambda_{i,v}/(N_v)^s),
\end{align*}
where the $\lambda_{i,v}$'s are the eigenvalues of $F_{v,\rho}$ (note
that the $\lambda_{i,v}$'s are algebraic numbers and hence may be
identified with complex numbers).  Put:
\[
L_\rho(s) = \prod_{v \notin S}\ \frac{1}{P_{v,\rho}((N_v)^{-s})}\,.
\]
This is a \emph{formal} Dirichlet series
$\sum\limits_{n=1}^{\infty} a_n/n^s$, with coefficients in $\Q$.  In
all known cases, there exists a constant $k$ such that
$|\lambda_{i,v}| \leq (N_v)^k$, and this implies that $L_\rho$ is
convergent in some half plane $\mathrm{R}(s) > C$; one conjectures it
extends to a meromorphic function in the whole plane.  When $\rho$
comes from \textit{l}-adic cohomology, there are some further
conjectures on the zeros and poles of $L_\rho$, cf.~Tate
\cite{tate-cycles}; these, as indicated by Tate, may be applied to get
equidistribution properties of the Frobenius elements, cf.~Appendix.

\begin{remarks}
\begin{nlist}
\item One can also associate \textit{L}-functions to
$E$-\emph{rational} systems of $\lambda$-adic representations
(\S\ref{sec:I.2.3}, Remark), where $E$ is a number field, once an
embedding of $E$ into $\C$ has been chosen.

\item We have given a definition of the local factors of $L_\rho$ only
at the places $v \notin S$.  One can give a more sophisticated
definition in which local factors are defined for all places, even
(with suitable hypotheses) for primes at infinity (gamma factors);
this is necessary when one wants to study functional equations.  We
don't go into this here.  $\longrightarrow$ [see \cite{deligne-valeurs},
\cite{serre-facteurs}.]

\item Let $\phi(s) = \sum\limits_{n=1}^{\infty} a_n/n^s$ be a
Dirichlet series.  Using the theorem in \S\ref{sec:I.2.3}, one sees
that there is (up to isomorphism) at most one semi-simple system
$\rho = (\rho_\ell)$ over $\Q$ such that $L_\rho = \phi$.  Whether
there does exist one (for a given $\phi$) is often a quite interesting
question.  For instance, is it so for Ramanujan's
$\phi(s) = \sum\limits_{n=1}^{\infty} \tau(n)/n^s$, where $\tau(n)$ is
defined by the identity
\[
x \prod_{n=1}^{\infty} (1 - x^n)^{24} = \sum_{n=1}^{\infty} \tau(n)x^n?
\]
There is considerable numerical evidence for this, based on the
congruence properties of $\tau$ (Swinnerton-Dyer, unpublished); of
course, such a $\rho$ would be of dimension $2$, and its exceptional
set $S$ would be empty.  $\longrightarrow$ [proved by Deligne: see
\cite{carayol}, \cite{deligne-formes}, \cite{langlands-modular},
\dots]

More generally, there seems to be a close connection between modular
forms, such as $\Sigma\tau(n)x^n$, and rational (or algebraic)
\textit{l}-adic representations; see for instance Shimura
\cite{shimura-recip} and Weil \cite{weil-dirichlet}.  $\longrightarrow$
[see also \cite{carayol}, \cite{deligne-valeurs},
\cite{langlands-modular}, \cite{langlands-automorphic}, \cite{ohta},
\cite{wiles}.]
\end{nlist}
\end{remarks}

\begin{examples}
\begin{enumerate}[label=\arabic*.,leftmargin=2.2em,itemsep=0.15ex,topsep=0.4ex]
\item If $G$ acts through a \emph{finite} group, $L_\rho$ is an Artin
(non abelian) \textit{L}-series, at least up to a finite number of
factors (cf.~\cite{artin-collected}).  All Artin \textit{L}-series are
gotten in this way, provided of course one uses $E$-\emph{rational}
representations (cf.~Remark 1) and not merely rational ones.

\item If $\rho$ is the system associated with an elliptic curve $E$
(cf.~\S\ref{sec:I.1.2}), the corresponding \textit{L}-function gives
the non-trivial part of the zeta function of $E$.  The symmetric
powers of $\rho$ give the zeta functions of the products
$E \times \ldots \times E$, cf.~Tate \cite{tate-cycles}.
\end{enumerate}
\end{examples}
